\begin{frame}{Learning Outcomes}
  By the end of this lecture you should be able to:
  \begin{itemize}\footnotesize\itemsep0.25em
    \item \textbf{Choose} an adaptation recipe for a pretrained classifier, weighing in-distribution accuracy against accuracy under distribution shift.
    \item \textbf{Explain} why fine-grained recognition is hard and how part-aware transformers, foundation features and taxonomic labels help.
    \item \textbf{Derive and compare} pair-, proxy- and margin-based metric-learning losses and the ranking metrics used to judge them.
    \item \textbf{Design} an image-retrieval pipeline, from global descriptor to re-ranking to approximate nearest-neighbour search.
    \item \textbf{Describe} the re-identification protocol and the methods behind current results.
    \item \textbf{Compute and compare} scores that flag inputs from unknown classes, and evaluate them.
    \item \textbf{Explain} how foundation models detect unknown classes zero-shot and discover new categories in unlabelled data.
  \end{itemize}
  \vspace{0.3em}
  {\footnotesize \textbf{Prerequisite knowledge:} softmax, $k$-NN, $k$-means, mean average precision; CNN and ViT backbones, transfer learning; contrastive learning (InfoNCE, NT-Xent), DINO, DINOv2, CLIP.\par}
\end{frame}
